% Invented fixture: beamer's {overprint} as alternation. One item inked per
% overlay step and the others absent — never several readings stacked, which
% is what keeping the body did. Three items, one an inline run, one a list,
% one two paragraphs, so an item holding block content steps whole; a second
% frame puts content before the first item, which stands on every step; a
% third frame holds the reading the rewrite refuses to number, kept once.
\documentclass[aspectratio=169]{beamer}

\begin{document}

\begin{frame}{Three readings}
\begin{overprint}
\onslide<1> The lone opening reading.
\onslide<2>
\begin{itemize}
\item A listed point on the middle step.
\item A second listed point beside it.
\end{itemize}
\onslide<3>
The closing reading opens a paragraph.

And ends on a second paragraph.
\end{overprint}
\end{frame}

\begin{frame}{A line that stands}
Standing above every reading.
\begin{overprint}
\onslide<1> Shown on the opening step only.
\onslide<2> Shown on the later step only.
\end{overprint}
\end{frame}

\begin{frame}{Unnumbered}
\begin{overprint}
\onslide A reading kept as written.
\end{overprint}
\end{frame}

\end{document}
